\documentclass[11pt,a4paper]{article}
\usepackage[margin=22mm]{geometry}
\usepackage{amsmath,amssymb,amsthm,booktabs,microtype}
\usepackage[T1]{fontenc}
\usepackage{lmodern}
\usepackage[colorlinks=true,linkcolor=blue,urlcolor=blue]{hyperref}
\hypersetup{pdftitle={Variable-threshold hole bounds and the audited bound 90},pdfauthor={Sum-of-Subset research project}}
\newcommand{\FS}{\operatorname{FS}}
\newcommand{\Z}{\mathbb Z}
\newcommand{\T}{\mathcal T}
\newtheorem{theorem}{Theorem}[section]
\newtheorem{lemma}[theorem]{Lemma}
\newcommand{\RootGroups}{27,802}
\newcommand{\RootSystems}{371,986}
\newcommand{\FinalSystems}{372,137}
\newcommand{\Vectors}{995}
\newcommand{\MinGap}{3/151}
\newcommand{\NegativeControls}{43}
\title{Subset Sums in a Fourfold Sumset\\\large Variable thresholds for persistent holes\\and the audited bound $N\le90$}
\author{Sum-of-Subset research project}
\date{18 October 2026}
\begin{document}
\maketitle
\begin{abstract}
We continue the independently audited bound $N\le91$ by excluding
cardinality 91. A variable threshold at each prefix improves the
persistent-hole argument: at size 37 and sufficiently many holes, its
defect floor rises from 144 to 193. A two-stage lattice comparison uses
the one-budget bound first and reserves the stronger hole witnesses for
the remaining range. This closes the inner lattice gate without a new
defect-sensitive span lemma. The inverse theorem is used only when both
supports are non-progressions; progression summands are handled by an
exact gap formula, including a two-residue exception. We then reconstruct
the complete necessary covering models at the new endpoint. Exact integer
certificates and a standard-library verifier accompany the report.
The stronger archived bound 84 is not asserted to have received a
complete independent replay.
\end{abstract}

\section{Result and the inherited reduction}
For a finite real set $A$, $\FS(A)$ denotes all subset sums, including
zero, and $4A=A+A+A+A$ allows repetitions. Let $N$ be the largest
cardinality for which $\FS(A)\subseteq4A$.
\begin{theorem}\label{main}
Every feasible set has at most 90 elements. Hence $71\le N\le90$.
\end{theorem}
The lower construction remains
$([-31,31]\cap\Z)\cup\{\pm93,\pm155,\pm217,\pm713\}$.
The previous audited bound is 91 \cite{previous}. Thus it suffices to
exclude size 91, and such a set must contain zero: otherwise adjoining
zero preserves feasibility and gives 92 elements. We use the corrected
integralization, covering, and peeling lemmas in \cite{foundation}.
No inheritance of feasibility by arbitrary subsets is assumed.

Write $T_r=r(r+1)/2$. For a positive $r$-set $X$, let
\[
 \mu(X)=|\FS(X)|-1-T_r,\qquad
 \rho(X)=\mu(X)+2-2\delta_0(X),
\]
where $\delta_0$ counts $x$ for which $2x$ is absent from the restricted
twofold sumset of $X\cup\{0\}$. We inherit $\rho\ge-8$ universally,
$\rho\ge0$ for $r\ge14$, and $\rho\le\mu$ because the maximum is
defective. Non-progression support has $\mu\ge r-3$; its defect is at
least $r-14$ for $r\ge6$ and at least $r-5$ for $r\ge17$ \cite{sharp}.
Let $f(r)$ denote the latter available floor.

Remove four original extremes on each side to obtain $D$, and perform
the two four-extreme peels $D\supset D_1\supset E$, where
$E=P\cup\{0\}\cup(-Q)$. At the present endpoint,
\[
 d=83,\quad e=75,\quad a+b=m=74,\quad I=(a-37)^2,
 \quad\Lambda=1+T_a+T_b=1407+I,\quad\Omega=516-I.
\]
The exact inherited budget is
\begin{equation}\label{budget}
 \rho+2\vartheta+c+\T=\Omega,\qquad
 c=2|2D|+18d+85-|\FS(D)|,\qquad
 \T=\kappa+\ell_1+\ell_2+g,
\end{equation}
where $\rho=\rho(P)+\rho(Q)$, $\mu=\mu(P)+\mu(Q)$, and
$\kappa=|\FS(P)-\FS(Q)|-|\FS(P)|-|\FS(Q)|+1$.
All displayed slacks other than $\rho$ are nonnegative.

The deficient-sign check is $1+T_{75}=2851>24\cdot91+77=2261$.
The paired bootstrap gives $q^2\le2648$ with even $q$, so $q\le50$
and $a,b\ge12$. At a side of size 12 or 13, respectively,
$\Omega=-109$ or $-60$, below the general lower bound $-8$ for the
combined defect. These sizes are impossible. Both sides are now at
least 14; nonnegative defect excludes size 14, where $\Omega=-13$.
Thus $15\le a\le59$, giving 45 oriented lattice cases.

\section{Variable thresholds for persistent holes}
We first prove the new general estimate. It uses only inherited
positive-side lemmas and the full-prefix central interval from
\cite{previous}, itself a direct application of Lev's Theorem 1
\cite{lev}. The latter states that an $i$-set in $[1,l]_{\Z}$ with
$l\le3i/2-2$ fills all subset sums between $2l-2i+1$ and its
complement about half the total. Consequently, if $Y$ has $i$ elements,
$\max Y\le i+\beta$, and $i\ge2\beta+4$, then
\begin{equation}\label{central}
 \FS(Y)=[0,S]_{\Z}\setminus(H\cup(S-H)),\quad
 H\subseteq[1,2\beta],\quad |H|\le\beta,\quad S=\sum Y.
\end{equation}
Indeed apply Lev with $l=i+\delta$, where $0\le\delta\le\beta$;
each missing lower sum is a missing coefficient in $[1,2\delta]$.

\begin{theorem}[A schedule of small-increment thresholds]\label{schedule}
Let $X$ be a positive integer $r$-set with hull-hole count
$z\le h_X<|\FS(X)|-1$. Choose integers $18\le L\le r$, $B\ge0$,
and $\beta_i\ge2$ for $L\le i\le r$. Suppose
\begin{align}
 2\beta_i&<z,& 2\beta_i+4&\le i,\label{local}\\
 B&\le T_{L-1}-r-1,&
 B&\le T_i+i-6\beta_i-r-7\quad(L\le i\le r),\label{escape}\\
 B&\le L-6+\sum_{i=L}^r(\beta_i-1).&&\label{accumulate}
\end{align}
Then $\rho(X)\ge B$.
\end{theorem}
\begin{proof}
Assume $\rho_r<B$. The strict upper hole condition forces $X$ to
be primitive; a gcd at least two would create at least
$|\FS(X)|-1$ holes. Its support is non-progression because $h_X>0$.
For increasing prefixes $X_i$, put
$y_i=|\FS(X_i)|-|\FS(X_{i-1})|-i\ge0$.
We have $\rho_i\ge\rho_{i-1}+y_i-2$, and $y_i=0$ at $i\ge6$
forces progression support.

If the last progression prefix has size $t\ge L-1$, its next
translate is disjoint: an on-lattice overlapping or abutting translate
would preserve progression support. Thus $y_{t+1}\ge T_t-t$.
All later prefixes are non-progressions, so their increments are at
least one. Starting with $\rho_t\ge0$ gives
$\rho_r\ge T_t-r-1\ge T_{L-1}-r-1\ge B$, a contradiction.
Every prefix from $L-1$ onward is therefore non-progression, and
$\rho_{L-1}\ge L-6$. If $y_i>\beta_i$ for every $i\ge L$,
the recurrence and \eqref{accumulate} already imply $\rho_r\ge B$.

Choose an $i\ge L$ with $1\le y_i\le\beta_i$.
The inherited small-increment inverse applies since
$i\ge2\beta_i+4\ge\beta_i+5$. It writes $X_i=tY$ with $Y$
primitive, $|Y|=i$, and $\max Y\le i+\beta_i$. Here $t$ is an integer.
If $t>1$, final primitivity forces a later first element off $t\Z$.
Its translate is disjoint, and the preceding calculation gives
$\rho_r\ge T_i-r-1\ge B$. Hence $t=1$ and \eqref{central} applies
with the fixed threshold $\beta_i$ chosen at this prefix.

Propagate its endpoint-hole description through later prefixes.
Suppose it holds at size $v\ge i$, with total $S$, and the next
element is $x$. If $x\ge S-4\beta_i$, the translate adds at least
$S-6\beta_i$ new points above $S$: for $x\le S$ the last $x$
positions lose at most $2\beta_i$ holes, while for $x>S$ the whole
translate is new. Using $\rho_v\ge v-5$ and later increments
at least one yields
\[
 \rho_r\ge S+v-6\beta_i-r-7
 \ge T_i+i-6\beta_i-r-7\ge B.
\]
Thus $x<S-4\beta_i$. Also $x>2\beta_i$ by increasing order.
The central intervals overlap or abut, the lower holes remain $H$,
and complement symmetry supplies the upper holes. Induction gives
$h_X\le2\beta_i<z$, the final contradiction.
\end{proof}
The threshold changes only in the search for the useful prefix.
Once that prefix is selected, its own $\beta_i$ remains fixed during
hole propagation. This distinction is essential.

For a proposed $B$ and $L$, the largest permitted local thresholds are
\begin{equation}\label{search}
 \beta_i(B)=\min\left\{
 \left\lfloor\frac{i-4}{2}\right\rfloor,
 \left\lfloor\frac{z-1}{2}\right\rfloor,
 \left\lfloor\frac{T_i+i-r-7-B}{6}\right\rfloor\right\}.
\end{equation}
The checker requires every threshold to be at least two and verifies
\eqref{local}--\eqref{accumulate} exactly. Feasibility is monotone
as $B$ decreases, so the optional witness search uses integer bisection.
The theorem needs any valid schedule, not optimality of that search.
The previous constant-threshold bound is retained if stronger.

At $(r,z)=(37,215)$, take $L=22$, $B=193$, and
\[
 (\beta_{22},\ldots,\beta_{37})=
 (6,9,10,10,11,11,12,12,13,13,14,14,15,15,16,16).
\]
The AP escape bound is exactly 193, the accumulation is 197, and
every local jump bound is at least 193. The previous bound was 144. The same schedule already works for $z\ge33$.
At $z=15,19$, respectively, the new floors are 126 and 152,
compared with 123 and 144. The helper
\path{verification/nonap_hole_rho_schedule.py} records all schedules,
checks 3,311 parameter queries (2,480 strictly improved), and retains
the previous finite central-interval audit. A separate deterministic
check computes actual supports and diagonal defects for 2,432 integer
sets, including 1,102 with the required strict hole condition. These finite checks supplement the general proof above.

\section{A complete progression-support repair}
The raw sufficient inverse margin is now
$J=T_{\min(a,b)}-2-\Omega$, with minimum $-49$.
We use Stanchescu's Corollary 2 \cite{stanchescu} only under the
sufficient strict condition
$\kappa<\min(|\FS(P)|,|\FS(Q)|)-3$.
If both supports are non-progressions, put $r=\min(a,b)$ and $t=74-r$.
The smaller-support margin is at least
\[
 J+(r-3)+f(r)+f(t)\ge36,
\]
and the larger baseline margin $T_t-2-\Omega$ is at least 185.
This gives a common core lattice, side holes $h_P,h_Q\le\kappa$,
and strict side-hole bounds $h_X<|\FS(X)|-1$.
Both positive sets are integral on this lattice.

If a support is a progression, we instead use an elementary exact formula.
For $W=\{w_1<\cdots<w_v\}\subset\Z$ and $J_u=[0,u-1]_{\Z}$,
\begin{equation}\label{gap}
 |W+J_u|=v+u-1+
 \sum_{i=1}^{v-1}\min(w_{i+1}-w_i-1,u-1).
\end{equation}
Add the translated intervals in increasing order to prove this.
For a progression of step $t$, apply the formula in each occupied
residue class. If there are $s$ classes, the excess is at least
$(s-1)(u-1)$.

We also need a subset-sum observation. For positive
$x_1<\cdots<x_b$, the largest consecutive gap in the support is
$\max_i(x_i-\sum_{j<i}x_j)$. Indeed adding $[0,L-1]$ fills one
interval precisely when $x_i\le L+\sum_{j<i}x_j$ for every $i$,
by induction; a failed test leaves a gap that later positive elements
cannot fill. If exactly one support gap has length at least $u$,
a prefix gap attaining the maximum and its complement reflection
must be the same gap. Their endpoints force $i=b$.
The last coefficient dominates the previous total, and therefore
\begin{equation}\label{dominant}
 \mu(X)\ge2(1+T_{b-1})-1-T_b=T_{b-2}.
\end{equation}

Call the progression $U=t[0,u-1]$, and the other support $V$.
If its side has at least 37 coefficients, then
$u-1\ge T_{37}=703>\Omega$. Formula \eqref{gap} forces $V$
into one residue class and forbids any gap of length at least $u$.
On that lattice $h_U=0$, $h_V=\kappa$, and $z=0$.

It remains that the progression side has $15\le a\le37$ coefficients,
and the other side has $b=74-a\ge37$. Exact quadratic arithmetic gives
\begin{equation}\label{apranges}
 \Omega<2T_a,\qquad \Omega-T_a\le47,\qquad T_a+1>2b.
\end{equation}
There are at most two residue classes in $V$.
In one class, there can be at most one gap of length at least $u$.
Such a gap would give \eqref{dominant} and
$\rho(Q)\ge T_{b-2}+2-2b\ge558>516$, impossible.
So this case also gives $h_U=0,h_V=\kappa,z=0$.

In two classes, one off-lattice coefficient would give two disjoint
copies of a $(b-1)$-element subset-sum support and the same contradiction.
There are at least two off-lattice coefficients. Their common nonzero
residue must have order two, since their sum is also represented.
Normalize the half-step so that $U=2[0,u-1]$ and $Q$ has at least
two odd coefficients. Let $w$ be its least odd coefficient, and let
$h_0,h_1$ be the two fiber hull-hole counts, measured with step two.
No fiber gap reaches $u$ in its own coordinates: the baseline excess
$u-1$ plus that gap would contradict \eqref{apranges}. Formula
\eqref{gap} and complement symmetry give
\begin{equation}\label{fibers}
 \kappa=u-1+h_0+h_1,\quad h_U=u-1,\quad
 h_Q=\kappa-u+w,\quad z=w-1.
\end{equation}
The $w-1$ term counts the integer positions outside the two fiber hulls.
These identities hold whether the total of $Q$ is even or odd.

We first show $w\le u$. If $w\ge u+1>2b$ and $s\ge2$ coefficients
are odd, distinctness gives
\[
 \sum Q\ge(b-s)(b-s+1)+sw+s(s-1).
\]
This quadratic is increasing for $s\ge2$. Using \eqref{fibers},
$|V|\ge b^2-3b+2u+6-\kappa$.
But $\mu_Q+\kappa\le\Omega+2b-2$, by \eqref{budget} and
$\rho_Q\ge\mu_Q+2-2b$, gives
$|V|\le T_b+\Omega+2b-1-\kappa$.
The lower minus upper margin is at least
$(b^2-11b)/2+2(T_a+1)+7-\Omega>0$ in every case.
Thus both side holes in \eqref{fibers} are at most $\kappa$.

In fact $w\le13$. For odd $w\ge15$, minimizing the same mass
quadratic over real $s$ gives
\[
 \sum Q-w-T_b\ge7b-289/8.
\]
On the other hand $\mu_Q+\kappa=\sum Q+u-w-T_b$ and
\eqref{apranges} bound this expression above by $2b+44$,
a contradiction for $b\ge37$. Therefore $z\le12$.

The progression side satisfies $R_P\le2(a+\mu_P)$ by subtracting
$2T_{a-1}$ from its total. The other side satisfies
$R_Q\le b+\mu_Q+\kappa-T_a-\mu_P+12$ by subtracting $T_{b-1}$.
Insert these bounds into the exact pre-lattice comparison, and use
$\mu\le\rho+2m-4$ and $\rho+\kappa+\ell_2\le\Omega$:
\[
 |\FS(D_1)|-2|\FS(E)|
 \le2543+2a-4T_a-2T_b\le-1157.
\]
The outer comparison has margin at least 569. Off-lattice points would
give disjoint support translates, so $D_1,D$ lie on the half-step lattice.
Only now may we use the sharp peeling bound
$\T\ge\kappa+156-3z\ge T_a+120$, contradicting
$\Omega-T_a\le47$. Its margin is at least 73.
This excludes the two-residue exception without a circular use of
post-lattice peeling.

The progression cases left have $z=0$. The general span lemma
$\max X\le |X|+\rho(X)+h_X+2$ for integer sets of size at least 15
requires no strict side-hole condition \cite{previous}. It gives inner
comparison $C+\Omega<0$ with $C=-945-I$, and the same outer comparison.
After lattice propagation, $\kappa\le\Omega-156$, which gives
$T_{\min(a,b)}-\kappa\ge109>0$. Strict side-hole bounds are now
available in these cases as well.
This progression repair rederives the analytic size-91 argument
in the archive \cite{archive}; none of its endpoint certificate files
is used here.

\section{A two-stage lattice comparison}
For the non-progression case the core is already normalized, so the
new hole theorem is available before the peeled layers are integral.
Write $s=|\FS(E)|$ and $z=h_P+h_Q-\kappa$, so that
$0\le z\le\min(h_P,h_Q)$ and both strict side-hole bounds hold.
The inherited inner comparisons are
\begin{equation}\label{inner}
 |\FS(D_1)|-2s\le
 \min\{C+\Omega+\rho-\mu+2z,\ C+3\Omega-\rho-\mu\},
 \quad C=-945-I.
\end{equation}
Since $\rho\le\mu$, the first is $-429-2I+2z<0$ whenever
$z\le214+I$. In the remaining range $z\ge215+I$, first use the
non-progression floors. Any unresolved orientation has $a,b\ge26$.
For each side choose the best inherited effective-hole witness
$\mu(X)\ge u_r=r+C_h+1$ with $2\beta<215+I$.
Its two strict hypotheses are checked again by the verifier.
Combine this with
\[
 \rho(X)\ge\max\{F(r,u_r),\mathcal R(r,215+I)\},
\]
where $F$ is the exact prefix inverse \cite{prefix} and
$\mathcal R$ is the maximum of the old and new persistent-hole floors.
All 45 oriented comparisons are strictly negative: the first-stage
maximum is $-1$, and the remaining second-stage maximum is $-13$.
At balance the previous constant-threshold bound left the comparison
at $+3$. Now $u_{37}=156$ and $\mathcal R(37,215)=193$, giving
$-945+3\cdot516-2(156+193)=-95$.

The outer margin is
$\Lambda-(4d-10+\Omega)=569+2I>0$.
Thus both peeled layers lie on the core lattice.
Hole propagation and sharp four-extreme peeling \cite{overlap} now give
\begin{equation}\label{peeling}
 w\le4z,\quad \T\ge\kappa+156-3z,\qquad
 \T\ge\max(52,z,156-2z).
\end{equation}
No state with a side of size 15 or 16 survives these budgets, so
all later inverse and hole-helper calls are within the range $r\ge17$.

\section{Covering models at the new endpoint}
We preserve the actual-set-to-model construction from \cite{previous}:
the same eighteen ordinary units, first actual contributor orientation,
split-pair choices, integer row-lattice witnesses, loss variables,
subgroup overlap inequalities, six core cutoff positions, and 85 trimmed
sumset rows. The possible empty-unit count is still at most four.
All 10,880 four-empty modes are retained, including the 30 modes with
only partial integer gap information. No rational division is used to
infer an integer coordinate.

The four low-hole families are
\begin{center}\small
\begin{tabular}{lrrrr}\toprule
Family&$z$&$\rho_P+\rho_Q\ge$&$\mu(X)\ge$&$c+I\le$\\\midrule
Z&0&0&0&360\\
P&1--2&64&$r-3$&300\\
Q&3--4&76&$r+6$&292\\
R&5--14&97&$2r-5$&291\\\bottomrule
\end{tabular}
\end{center}
The last three caps equal $360-\rho_{
\mathrm{floor}}+2z_{
\mathrm{endpoint}}$.
The effective-hole rows and exact prefix inverse certify the Q and R
defect floors for $17\le r\le57$. Each positive family transports
$z$ to its own endpoint $Z$ and adds $Z-z$ to both mass coordinates.
Every source row is checked for the required nonpositive change;
the auxiliary masses need not describe another actual set.
For empty count $j$, the low parameter interval is exactly the set of
integers $a$ satisfying $(a-37)^2\le C_{
\mathrm{family}}-83j$.
These envelopes contain 369,640 integer low roots, including any
already impossible scalar states; enlarging the necessary model domain
is harmless.

For $z\ge15$ the scalar enumeration leaves 1,320 states; $c\le140<2d$,
so only $j=0,1$ occur. The all-active comparison leaves 196 states
at $j=0$ and the empty-unit budget leaves 66 at $j=1$.
They form 42 groups with fixed side sizes and support/defect floors.
For comparison, the same post-lattice scalar relaxation with only the
previous constant-threshold floor has 3,596 states, cap 176, and 56
two-empty states; the variable thresholds reduce these to 1,320, 140,
and zero. This comparison alone is not an endpoint exclusion.
The same real variable $z$ is retained across each group's full interval.
With $E_*=\Omega-\rho_p-\rho_q$ and $\delta=c-83j$, use all four caps
\[
 \delta\le c_{\max}-83j,\quad
 \delta+z\le E_*-83j,\quad
 \delta\le E_*-52-83j,\quad
 \delta-2z\le E_*-156-83j.
\]
The verifier checks 109,514 coarse-source and 211,432 nested-source
row implications, including every new full-translate row.
It does not combine unrelated maxima of $c$ and $z$, nor choose the
reflected all-active orientation independently of the actual contributor.

For clarity, all dimension-dependent replacements in the inherited
covering model are listed below. Put $K=|2D|$, let $R_s,C_s$ be the
largest and second largest core magnitudes, and let $M_s$ be the
largest magnitude in $D$.
\begin{center}\small
\begin{tabular}{lr}\toprule
Quantity&New value\\\midrule
Ordinary unit capacity $d$; core size $e$&83; 75\\
Global quantile pair; core quantile pair&83/82; 79/78\\
Ordinary quantile pair&opposite size $+5$ / own size $+4$\\
Basic leading threshold $2d-1$&165\\
Nested leading threshold $2d-9$&157\\
Second-core leading threshold $2d-11$&155\\
Core loss threshold $d-4$&79\\
Span budget $R_p+R_q+\delta-4z$&$438-I-83j$\\
Cardinality budget $S_p+S_q+\delta-2K-4z$&$1578-83j$\\\bottomrule
\end{tabular}
\end{center}
Each threshold is applied to its own proved residual cap and charged
leading side. The trimmed increment remains $83s-T_{s-1}$ for removing
$s$ extreme elements. The coarse bound remains
$K\le M_p+M_q+R_p+R_q+21$.
For high groups, $140<155$ licenses the deeper leading quantiles;
after one empty unit, $140-83=57<79$ licenses core quantiles.
The actual selected cap, not these global maxima, determines each row.

The coarse model is tried first. When needed, the nested model and
the inherited exhaustive disjunctions are used: density, ordinary versus
outside-center quantiles, core-loss alternatives, six inner positions,
and both-or-sole contributors for a split pair. Every binary child must
be present. No solver-selected cut is accepted unless reconstructed from
one of these proved alternatives.

\subsection{Using the location of an empty own translate}
There is one additional low-hole refinement. At $z=0$, the target is
the full interval $[1,S_p]_{\Z}$ on the positive side. Suppose an
integer own center $u=c_p+p_{i+1}$ is empty, but a lower own center
$v=c_p+p_{h+1}$, $h>i$, contributes. Then $u>v$ and $u>S_p$.
Indeed $0<u\le S_p$ contributes through $u+0$. If $u\le0$, choose
$d'\in D$ with $1\le v+d'\le S_p$. Then
$1\le u+d'\le d'\le M_p\le S_p$, also a contribution.
Thus the new necessary row is
\[
 -u+S_p\le-1.
\]
Reflection gives the negative-side row. We activate it only for
zero-hole groups with a complete integer-center basis, an explicitly
missing own single, and a retained lower own single.

When $i=0$, the four original fixed labels
$C^+_{1j}=c_p+p_1+p_j$ ($1\le j\le4$) all exceed $S_p$.
They therefore contribute four additional fixed losses. The original
fixed cover has 110 labels and 26 disjoint groups, each supplying a
nonpositive label. These four high labels are positive, so they cannot
be any of those nonpositive deletions; their losses are additional.
The reflected four labels are separate labels. Counting labels rather
than distinct values only enlarges the covering upper bound and is safe.

Let $f$ be four times the number of sides meeting this largest-empty
condition. Then both the ordinary loss sum and the combined subgroup
loss sum are at most $\delta-f$. The same correction is made in the
sole-contributor branch. It is one loss charged consistently in alternative
necessary inequalities, not a sum of repeated gains. These deductions
also rederive the applicable fixed-label argument in \cite{archive}.
The verifier reconstructs the original 110-label/26-group cover and
checks 10,430 small configurations, including 8,739 high-center tests
and 22,952 high fixed-label tests. The arbitrary-size proof is the
interval argument just given.

\subsection{Full translates, a common deficit, and clipping}
The nested model also uses all of $D$ in selected overlap comparisons,
rederiving the full-translate argument in \cite{archive}. On the
positive side let the leading component be $c_p+2D$, and select own
components $c_p+p_i+D$ and cross components $c_p-q_j+D$.
All these centers are integral in the modes where these rows are used.
If $k$ own and $l$ cross components are selected, their total nominal
capacity is $C=K+(k+l)d$. Let $p_*,q_*$ be their largest selected
outer magnitudes when present. The union lies in $[L,U]_{\Z}$, with
\[
 L=c_p-\begin{cases}q_*+M_q,&l>0,\\2M_q,&l=0,\end{cases}
 \qquad
 U=c_p+\begin{cases}p_*+M_p,&k>0,\\2M_p,&k=0.\end{cases}
\]
For the actual target $F_p=\FS(D)\cap(0,\infty)$ and subgroup union
$W$, put $G=C-|W\cap F_p|$. Every member loss $\lambda_v$ satisfies
\[
 G\ge C-(U-L+1)+\lambda_v,\qquad G\ge\sum_v\lambda_v.
\]
Indeed $C-|W\cap F_p|=(C-|W|)+|W\setminus F_p|$, and the outside
union contains each member's outside set. Nonemptiness of a selected
ordinary component further gives the clipped bounds
\[
 G\ge C-U,\qquad G\ge C-S_p+L-1,\qquad G\ge C-S_p.
\]
These are bounds for the actual target deficit; no selected member loss
is added to a clipped bound.

For each side, the larger family consists of the leading component,
all retained own components, and retained unpaired cross components
with labels 6 and 7 (outer indices 3 and 4). Its variable $\Gamma_p$
is the actual capacity deficit of the whole family. For any selected
subgroup, $\Gamma_p$ dominates its deficit plus the individual losses
of members outside that subgroup. Taking suffixes of the own and cross
lists gives all implemented mixed rows. The two larger families use
disjoint component labels, so
\[
 \Gamma_p+\Gamma_q+\sum_{\text{unmatched }v}\lambda_v\le\delta-f.
\]
The earlier block inequalities remain separate necessary inequalities;
their gains are not added to this new allocation. All-subset own-only
rows analogously strengthen the earlier own-block variables. The
49-set sanity audit checks 795 subgroup inequalities. Positive-hole
mass transport preserves these rows because each mass coefficient is
zero or negative.

\subsection{Exact fixed-label indicators in the four-empty exception}
Only the zero-hole four-empty mode with missing units
\texttt{pair1,pair2,p0,q0} enables the fixed-score extension, in either
structural branch. It already charges $f=8$. For each of the other
102 original fixed labels $v$, introduce an auxiliary $H_v\in[0,1]$,
realized as the indicator that the score exceeds its own target mass.
For all 110 labels introduce $N_v\in[0,1]$, realized as the indicator
that the score is nonpositive. Each of the 26 disjoint negative groups
requires the sum of its $N_v$ to be at least one. Thus
\[
 F_{\rm extra}=\sum_{v\text{ not already charged}}H_v+
               \sum_{v}N_v-26\ge0
\]
is an additional actual fixed loss. The eight already charged labels
are omitted from the high sum. High and nonpositive events are disjoint,
so this loss may be added to each complete deficit allocation.

All score forms are integer combinations of the eight proved integral
outer coordinates, and both masses are integral in these zero-hole
groups. For a score $f_v$ with target mass $S$, the exhaustive splits are
\[
 (f_v\le S,\ H_v=0)\quad\text{or}\quad
 (f_v\ge S+1,\ H_v=1),
\]
\[
 (f_v\le0,\ N_v=1)\quad\text{or}\quad
 (f_v\ge1,\ N_v=0).
\]
Their actual indicator realizations justify the necessary relaxations
and both branches. Universally positive score forms, recognized by
nonnegative coefficients in positive outer-gap coordinates, have
$N_v=0$ directly. The eight charged high labels also have $N_v=0$.
The helper \path{verification/fixed_scores.py} regenerates every form
from the covering definitions, transforms each of the 26 group sums
to gap coordinates, checks its nonpositive coefficients, and checks
disjointness. Stored coefficient vectors alone are not trusted.

\section{Exact certification and reproduction}
The certificate bundle contains \Vectors{} distinct nonnegative integer
multiplier vectors. It covers \RootGroups{} structural groups,
\RootSystems{} integer roots, and \FinalSystems{} final leaf systems.
Of the 27,802 groups, 27,250 close in the coarse model and 552 use
the nested model. There are 27,836 symbolic interval checks, 80,170
integer extremum evaluations, and 25 tree roots with 151 binary nodes
and 176 tree leaves. The maximum tree depth is 29. The splits comprise
18 core-loss, 13 inner-position, 13 high-score, and 107 nonpositive-score
alternatives. The minimum normalized strict contradiction gap is $\MinGap$.
For every leaf with variables $x\ge0$ and necessary rows $Ax\le b$,
the checker verifies $w\ge0$, $wA\ge0$, and $wb<0$ using integers
and exact rational arithmetic. The generator may use SciPy to discover
$w$; no floating-point result is accepted as a proof.

The checker reconstructs the whole root universe, each allowed
disjunction, and all parameter intervals. Quadratic right sides are
maximized at endpoints and the nearest integers to any interior vertex.
The exhaustive option additionally instantiates every integer parameter
and compares every row with its symbolic formula. The high-hole rows
cover real intervals and hence all actual integer hole counts.
Negative controls reject \NegativeControls{} deliberately malformed
proof inputs, including missing roots, omitted children, invalid
integer division, nonstrict schedule thresholds, and incomplete schedules.

From the repository root:
\begin{verbatim}
python -I -S -B verification/upper90.py --output .work/upper90.json
python -I -S -B verification/upper90.py --exhaustive \
  --output .work/upper90-exhaustive.json
python experiments/generate_upper90_certificates.py \
  --output .work/upper90-regenerated.json
\end{verbatim}
Only the optional generator needs NumPy and SciPy. The Windows entry
point is \path{verify_upper90_windows.bat}. Complete run receipts,
source hashes, scalar witnesses, and compressed full output are under
\path{results/2026-10-18/upper90/}. Ordinary and optimized-Python runs,
full integer expansion, independent certificate regeneration, and the
unchanged upper-91 baseline are recorded separately.

The result is a continuation of the audited chain, not a claim that
the archived bound 84 is false or has been completely rechecked.
The arbitrary-set mathematical lemmas remain part of the proof;
the program is an exact finite certificate checker, not a proof-assistant
formalization of the entire argument.

\begingroup\small
\begin{thebibliography}{9}\setlength{\itemsep}{1pt}
\bibitem{previous} Sum-of-Subset project,
\emph{Anchored hole bounds and integer center bases, and the audited bound
$N\le91$}, \path{reports/upper91_from92_20261017.tex}.
\bibitem{foundation} Sum-of-Subset project,
\emph{From the audited bound 115 to 103},
\path{reports/upper103_from115_20261014.tex}.
\bibitem{sharp} Sum-of-Subset project,
\emph{Sharp non-progression defects and the audited bound 94},
\path{reports/upper94_from95_20261016.tex}.
\bibitem{prefix} Sum-of-Subset project,
\emph{An integer prefix inverse and the audited bound 93},
\path{reports/upper93_from94_20261016.tex}.
\bibitem{overlap} Sum-of-Subset project,
\emph{Sharp peeling and joint covering losses},
\path{reports/upper95_from96_20261016.tex}.
\bibitem{archive} Sum-of-Subset project,
\emph{Complete report for the archived upper bound 86}, section
``A common core lattice at size 91'',
\path{reports/upper86_complete_20261005.tex}.
\bibitem{lev} V.~F. Lev, \emph{On consecutive subset sums},
Discrete Mathematics 187 (1998), 151--160, Theorem 1.
\href{https://doi.org/10.1016/S0012-365X(98)80006-X}{doi:10.1016/S0012-365X(98)80006-X}.
\bibitem{stanchescu} Y. Stanchescu,
\emph{On addition of two distinct sets of integers},
Acta Arithmetica 75 (1996), 191--194, Corollary 2.
\href{https://doi.org/10.4064/aa-75-2-191-194}{doi:10.4064/aa-75-2-191-194}.
\end{thebibliography}
\endgroup
\end{document}
